\exercisesfor{1}

\begin{enumerate}
\def\labelenumi{\arabic{enumi}.}
\item
  Suppose that K is a field of characteristic 0. Let E be a cocommutative bigebra of finite rank over K. Show that \(P(E) = \{0\}\) (apply Theorem 1 of no. 6).
\item
  Let E be a bigebra such that \(\mathrm{P}(\mathrm{E})=\{0\}\) and let \((\mathrm{E}_{n})_{n\geqslant0}\) be a filtration of E compatible with its bigebra structure. Show by induction on n that \(E_{n}^{+}=\{0\}\) for all \(n\geqslant0\) and deduce that \(E^{+}=\{0\}\) , i.e.~that E reduces to K.
\item
  Let \(G\) be a monoid and \(E = K[G]\) its algebra (Algebra, Chapter III, §2, no. 6).
\end{enumerate}

\begin{enumerate}
\def\labelenumi{(\alph{enumi})}
\item
  Show that there exists on E one and only one cogebra structure such that \(c(g) = g \otimes g\) for all \(g \in G\) ; this structure is compatible with the algebra structure on E and makes E into a cocommutative bigebra whose counit \(\varepsilon\) is such that \(\varepsilon(g) = 1\) for all \(g \in G\) .
\item
  Show that every primitive element of E is zero. Deduce that, if \(G \neq \{e\}\) , the bigebra E admits no filtration compatible with its bigebra structure (apply Exercise 2).
\item
  Suppose that \(\mathbf{K}\) is an integral domain. Show that the elements of \(\mathrm{G}\) are the only non-zero elements \(x\in \mathrm{E}\) such that \(c(x) = x\otimes x\) . Show that \(\mathrm{E}\) has an inversion \(i\) (cf.~Algebra, Chapter III, § 11, Exercise 4) if and only if \(\mathrm{G}\) is a group and that in that case \(i(g) = g^{-1}\) for all \(g\in \mathrm{G}\) .
\end{enumerate}

\begin{enumerate}
\def\labelenumi{\arabic{enumi}.}
\setcounter{enumi}{3}
\item
  Let \(\mathbf{E}\) be a cocommutative bigebra and let \(\mathbf{G}\) be the set of \(g \in \mathbf{E}\) such that \(\varepsilon(g) = 1\) and \(c(g) = g \otimes g\) . Show that \(\mathbf{G}\) is stable under multiplication and that it is a group if \(\mathbf{E}\) has an inversion. Show that, if \(\mathbf{K}\) is a field, the elements of \(\mathbf{G}\) are linearly independent over \(\mathbf{K}\) .
\item
  Let E be a bigebra and let \(E' = \text{Hom}(E, K)\) be its dual with the algebra structure derived by duality from the cogebra structure on E (Algebra, Chapter III, §11, no. 1). Let m be the kernel of the homomorphism \(u \mapsto u(1)\) of \(\mathbf{E}'\) onto \(\mathbf{K}\) ; it is an ideal of \(\mathbf{E}'\) . Show that, if \(u \in \mathfrak{m}^2\) and \(x \in \mathrm{P}(\mathrm{E})\) , then \(u(x) = 0\) . When \(\mathbf{K}\) is a field show that \(\mathrm{P}(\mathrm{E})\) is the orthogonal of \(\mathfrak{m}^2\) in \(\mathbf{E}^+\) ; deduce that \(\dim \mathrm{P}(\mathrm{E}) \leqslant \dim \mathfrak{m} / \mathfrak{m}^2\) and that we have equality if \(\mathbf{E}\) is of finite rank over \(\mathbf{K}\) .
\item
  Let E be a bigebra and let \(u \in \mathrm{E}' = \operatorname{Hom}(\mathrm{E}, \mathrm{K})\) (cf.~Exercise 5). Let \(g \in \mathrm{E}\) be such that \(\varepsilon(g) = 1\) and \(c(g) = g \otimes g\) . Show that the mapping \(u \mapsto u(g)\) is an algebra homomorphism of \(\mathrm{E}'\) into \(\mathrm{K}\) and that the mapping \(y \mapsto gy\) is an endomorphism of the cogebra \(\mathrm{E}\) ; show that this endomorphism is an automorphism if \(\mathrm{E}\) has an inversion.
\item
  Suppose that K is a field. Let E be a bigebra satisfying the following conditions:
\end{enumerate}

\begin{enumerate}
\def\labelenumi{(\roman{enumi})}
\item
  E is cocommutative;
\item
  E has an inversion (Algebra, Chapter III, § 11, Exercise 4);
\item
  \(\mathrm{P(E)} = \{0\}\) ;
\item
  E is of finite rank over K.
\end{enumerate}

Let \(\mathbf{E}'\) denote the dual K-algebra of the cogebra E. Let G denote the set of elements \(g \in \mathbf{E}\) such that \(\varepsilon(g) = 1\) and \(c(g) = g \otimes g\) ; it is a group (cf.~Exercise 4).

\begin{enumerate}
\def\labelenumi{(\alph{enumi})}
\item
  Let \(g \in \mathbf{G}\) and let \(\mathfrak{m}_g\) be the kernel of the homomorphism \(u \mapsto u(g)\) of \(\mathbf{E}'\) into \(\mathbf{K}\) . Show that \(\mathfrak{m}_g = \mathfrak{m}_g^2\) (reduce it to the case \(g = 1\) using Exercise 6; then apply Exercise 5).
\item
  Suppose that K is algebraically closed. Show that the ideals \(m_{g}\) are the only maximal ideals of \(E'\) ; deduce that their intersection is \(\{0\}\) , that \(E'\) is identified with the product \(K^{G}\) and that E is identified with the bigebra K[G] of the finite group G (cf.~Exercise 3).
\end{enumerate}

\begin{enumerate}
\def\labelenumi{\arabic{enumi}.}
\setcounter{enumi}{7}
\item
  Show that the enveloping bigebra of a Lie algebra \(\mathfrak{g}\) has an inversion \(i\) and that \(i(x) = -x\) for all \(x \in \sigma(\mathfrak{g})\) .
\item
  Suppose that K is a Q-algebra. Let g be a Lie algebra admitting a K-basis and let U be its enveloping bigebra with the canonical filtration \((\mathrm{U}_{n})_{n\geqslant0}\) . Let \(x\in U^{+}\) and let n be an integer \(\geqslant1\) . Show that x belongs to \(U_{n}^{+}\) if and only if \(c^{+}(x)\) belongs to \(\sum_{\substack{i+j=n\\ i,j\geqslant1}}\mathrm{Im}(\mathrm{U}_{i}^{+}\otimes\mathrm{U}_{j}^{+})\) .
\item
  Suppose that K is a Q-algebra. Let E be a cocommutative K-bigebra with a filtration compatible with its bigebra structure. Show that the morphism
\end{enumerate}

\[
f _ {\mathbf {E}} \colon \mathrm{U} (\mathrm{P} (\mathrm{E})) \rightarrow \mathrm{E}
\]

defined in no. 4, Proposition 8 is an isomorphism, if P(E) is assumed to be a free K-module (same proof as for Theorem 1).

¶ 11. Let E be a bigebra and let I be a finite set. We are given a basis \((e_{\alpha})\)

of E indexed by the elements \(\alpha\) of \(\mathbf{N}^{\mathrm{I}}\) and such that:

\begin{enumerate}
\def\labelenumi{(\roman{enumi})}
\item
  \(e_{0} = 1\) ;
\item
  \(c(e_{\gamma}) = \sum_{\alpha + \beta = \gamma} e_{\alpha} \otimes e_{\beta}\) for all \(\gamma \in \mathbf{N}^{\mathrm{I}}\) .
\end{enumerate}

Condition (ii) implies that E is isomorphic as a cogebra to the cogebra \(\mathrm{TS}(\mathrm{K}^{\mathrm{I}})\) of symmetric tensors of \(\mathrm{K}^{\mathrm{I}}\) (cf.~Algebra, Chapter IV, § 5, no. 7).

\begin{enumerate}
\def\labelenumi{(\alph{enumi})}
\tightlist
\item
  Let \(\mathrm{E}' = \mathrm{Hom}(\mathrm{E},\mathrm{K})\) be the dual algebra of \(\mathbf{E}\) and let \(\Lambda = \mathrm{K}[[(\mathrm{X}_i)_{\mathrm{i}\in I}]]\) be the algebra of formal power series in indeterminates \(\mathbf{X}_i\) indexed by I. For all \(u\in \mathrm{E}'\) let \(\phi_u\) be the formal power series
\end{enumerate}

\[
\sum_ {\alpha} u (e _ {\alpha}) x ^ {\alpha},
\]

where \(x^{\alpha} = \prod_{i\in I}\mathrm{X}_{i}^{\alpha (i)}\)

Show that \(u \mapsto \phi_u\) is an isomorphism of \(E'\) onto \(\Lambda\) and that this isomorphism maps the topology of simple convergence on \(E'\) to the product topology on \(\Lambda\) .

\begin{enumerate}
\def\labelenumi{(\alph{enumi})}
\setcounter{enumi}{1}
\tightlist
\item
  Let \(c_{\alpha \beta \gamma}\) be the constants of structure of the algebra \(E\) relative to the basis \((e_{\alpha})\) . Then
\end{enumerate}

\[
e _ {\alpha} e _ {\beta} = \sum_ {\gamma} c _ {\alpha \beta \gamma} e _ {\gamma}.
\]

We write x instead of \((\mathbf{X}_{i})_{i\in\mathbb{I}}\) and y instead of \((\mathbf{Y}_{i})_{i\in\mathbb{I}}\) , where the \(Y_{i}\) are new indeterminates. Show that there exists a family \(f(\boldsymbol{x},\boldsymbol{y})=(f_{i}(\boldsymbol{x},\boldsymbol{y}))_{i\in\mathbb{I}}\) of formal power series in the variables x, y such that

\[
\boldsymbol {f} (\boldsymbol {x}, \boldsymbol {y}) ^ {\gamma} = \sum_ {\alpha , \beta} c _ {\alpha \beta \gamma} \boldsymbol {x} ^ {\alpha} \boldsymbol {y} ^ {\beta} \quad \text { for   all } \gamma \in \mathbf {N} ^ {\mathrm{I}},
\]

where we write, as above, \(x^{\alpha} = \prod_{i} X_{i}^{\alpha(i)}\) and similarly for \(y^{\beta}\) and \(f(x, y)^{\gamma}\) .

Show that \(f(x, y)\) is a formal group law over \(K\) of dimension \(n = \text{Card}(I)\) , in the sense of Chapter I, § 1, Exercise 24†; define a Lie algebra isomorphism of this formal group law onto the Lie algebra \(P(E)\) , which has basis the \(e_{\alpha}\) with \(|\alpha| = 1\) .

\begin{enumerate}
\def\labelenumi{(\alph{enumi})}
\setcounter{enumi}{2}
\item
  Conversely, show that every formal group law over K in the x and y can be obtained by the above procedure, which is unique up to isomorphism.
\item
  When K is a Q-algebra, apply the above to the enveloping bigebra of a Lie algebra g with a finite basis over K. Deduce the existence (and uniqueness, up to isomorphism) of a formal group law with g as Lie algebra.
\item
  Give explicitly the bigebra corresponding to the formal group law in one parameter \(f(x, y) = x + y\) (resp. \(f(x, y) = x + y + xy\) ).
\end{enumerate}

¶ 12. Suppose that K is a field of characteristic p > 0.

\begin{enumerate}
\def\labelenumi{(\alph{enumi})}
\tightlist
\item
  Let \(\mathfrak{g}\) be a Lie \(p\) -algebra (Chapter I, § 1, Exercise 20) and let \(\tilde{\mathbf{U}}\) be its
\end{enumerate}


restricted enveloping bigebra (Chapter I, §2, Exercise 6). Show that there exists on \(\tilde{\mathbf{U}}\) one and only one bigebra structure which is compatible with its algebra structure and for which the elements of \(g\) are primitive. Show that \(\tilde{\mathbf{U}}\) is cocommutative and that \(\mathrm{P}(\tilde{\mathbf{U}}) = g\) .

\begin{enumerate}
\def\labelenumi{(\alph{enumi})}
\setcounter{enumi}{1}
\item
  If \(n\) is an integer \(\geqslant 0\) , let \(\tilde{\mathbf{U}}_n\) denote the vector subspace of \(\tilde{\mathbf{U}}\) generated by the products \(x_1\cdots x_n\) where \(x_i \in \mathfrak{g}\) for all \(i\) . Show that \((\tilde{\mathbf{U}}_n)_{n \geqslant 0}\) is a filtration of \(\tilde{\mathbf{U}}\) compatible with its bigebra structure.
\item
  Let \(\mathbf{E}\) be a bigebra and let \(\mathfrak{g} = \mathrm{P}(\mathbf{E})\) be the Lie algebra of its primitive elements. The mapping \(x\mapsto x^p\) leaves \(\mathfrak{g}\) stable and gives \(\mathfrak{g}\) a Lie \(p\) -algebra structure. Show that the canonical injection \(\mathfrak{g}\to \mathbf{E}\) can be extended uniquely to a bigebra morphism \(\tilde{\mathbf{U}}\rightarrow \mathbf{E}\) and that this morphism is injective (use the fact that, if \((x_{i})_{i\in I}\) is a basis of \(\mathfrak{g}\) with a total ordering, the monomials \(\prod_{i\in I}x_i^{n_i}(0\leqslant n_i < p)\) form a basis of \(\mathbf{U}\) (loc. cit.) and argue as in the proof of Lemma 2).
\end{enumerate}
% semantic-fragments: legacy-input-seam
\relax{}
